\section{ Knots as Mathematical Objects
}



\subsection{From Rope to Curve}
A mathematical knot is a closed curve in three-dimensional
space that has no self-intersections. Formally, it may be
represented by a continuous one-to-one map
\begin{equation}
    \mathscr{K} : S^1 \longrightarrow \mathbb{R}^3
\end{equation}
% \mathscr{H}
% \mathcal{\varkappa} : S^1 \to \R^3
%  \mathcal{K} : S^1 \to \R^3
where $S^1$ denotes the unit circle. Two knots are considered equivalent if one can be continuously deformed into
the other without cutting the curve or allowing it to pass
through itself \cite{adams2004}.

The simplest knot is the \underline{unknot}, which can be deformed into an ordinary circle. A nontrivial example is
the trefoil knot shown in Figure \ref{fig:knot}.

\begin{figure}[H]
    \centering
    \includegraphics[width=\columnwidth, angle=180]{template/trefoil_knot.png}
    \caption{ A planar drawing of a trefoil knot}
    \label{fig:knot}
\end{figure}


\subsection{Diagrams and Crossings
}
A knot diagram is a projection of a spatial knot onto a
plane. At every crossing, a small break indicates which
strand passes underneath. The main diagrammatic objects are:
% \begin{table}[h]
%     \centering
%     \resizebox{\columnwidth}{!}{%
%     \begin{tabular}{ll}
%         \texttt{Arc} & A continuous segment between two under crossings. \\
%          \texttt{Crossing} & A point where one projected strand passes
% over another. \\
%          \texttt{Orientation} &  A selected direction of travel around the
% knot.  \\
%     \end{tabular}
%     }
%     % \caption{Caption}
%     % \label{tab:placeholder}
% \end{table}

\begin{labeling}{Orientation}
\item[\texttt{Arc}]
A continuous segment between two under-crossings.

\item[\texttt{Crossing}]
A point where one projected strand passes over another.

\item[\texttt{Orientation}]
A selected direction of travel around the knot.
\end{labeling}

\cblue{A diagram may contain many crossings even when
the underlying knot is simple}, while \red{ an invariant must
remain unchanged under every legal deformation of the
knot.}


